\section{The complete physical incidence--inverse-gradient multiplier}
\label{sec:sm-singular-multiplier}
We prove a uniform integrability estimate for the actual multiplier
$I_m/|\nabla_qF|$ on the entire physical source.  The fixed-margin
assumption in the previous singular-volume payment is absent.
The mechanism uses the already proved endpoint-momentum kernel on
two different blocks of the same orbit.  Their intersection gives
more information than two separate marginal tails; no product law
or conditional independence is assumed.

Throughout $m\ge1$, $R\in[R_-,R_+]$, and
\[
 \zeta=(p_0,p_m),\qquad r=|\zeta|,\qquad
 Z_m=\frac{I_m}{|\nabla_qF|}=\frac{I_m}{Rr},\qquad
 I_m=1+\sum_{i=1}^{m-1}\frac{q_0^{d_i}}{c_i}.
\]
The identity $|\nabla_qF|=Rr$ is the first variation of the
actual reduced endpoint action, where $q=(\alpha_0,\alpha_m)$.
No clock, Euclidean norm or radius factor is changed.
Values on $r=0$ and on undefined singular collision records can
be assigned arbitrarily on their collision-area null sets; every
nonnegative integral below is equivalently defined by truncating
on regular records with $r>0$.  This convention concerns bulk
integrals and assigns no value to a separate singular paired trace.

Let $C_{\rm mom}$ be the complete kernel constant of
Theorem~\ref{thm:mc-kernel}, with its already proved physical curve
class, growth constant and all collision words.  Set
\begin{equation}\label{eq:sm-constants}
 A_0=\frac{C_{\rm mom}}4,\qquad
 A_1=\frac{C_{\rm mom}}\pi,\qquad
 K_\times=5A_0^{1/3}A_1^{2/3},\qquad
 J_\times=\frac1{R_-}
       \left(2\sqrt{A_0}+\frac{188}{3}K_\times^{3/4}\right).
\end{equation}
These constants have no dependence on $m$, an incidence margin,
an exact label, the target roof, or any outer microscopic width.

\begin{lemma}[Intersecting two actual momentum kernels]
\label{lem:sm-intersection}
For every $s,u>0$ and $1\le i<m$, the unconditioned physical
collision measure satisfies
\begin{equation}\label{eq:sm-intersection}
 \nu\{r<s\}\le A_0s^2,\qquad
 \nu\{r<s,\ c_i<u\}\le A_1s u^2.
\end{equation}
Both statements sum the complete actual source before taking any
supremum in radius, collision count or the index $i$.
\end{lemma}
\begin{proof}
Integrating \eqref{eq:mc-kernel} in the actual initial momentum
$p_0$ with $d\nu=d\alpha_0\,dp_0/(4\pi)$ gives
\[
 (p_0,p_l)_*\nu\le\frac{C_{\rm mom}}{4\pi}\,da\,dv
 \qquad\hbox{for every integer }l\ge1.
\]
This is the unconditioned version of \eqref{eq:mc-joint}; no
section factor belongs to it.  Use $l=m$ and dominate the area
of $\{(a,v):a^2+v^2<s^2\}$ by $\pi s^2$.  This proves the
first inequality, for all positive $s$ including $s>1$.

For the second inequality, retain the actual block from collision
$0$ to collision $i$.  Since $\{r<s\}\subset\{|p_0|<s\}$,
\[
 \nu\{r<s,c_i<u\}
 \le\nu\{|p_0|<s,\sqrt{1-p_i^2}<u\}.
\]
The $p_0$ interval has length at most $2s$.  If $0<u\le1$,
the second set of momenta has length
$2(1-\sqrt{1-u^2})\le2u^2$.  If $u>1$ its length is two,
which is again at most $2u^2$.  Apply the same kernel with $l=i$.
The product of these two coordinate lengths is only a Lebesgue
area under a proved joint-density upper bound.  It is not a
factorization of two physical probabilities.  The bound is
$C_{\rm mom}s u^2/\pi$, as claimed.  Future exact-record
tests can be dropped only on this positive upper side; all
collisions through $m$ remain in the source being bounded.
\end{proof}

\begin{lemma}[A mixed weak tail at every actual interior contact]
\label{lem:sm-mixed-tail}
For each $1\le i<m$ and every $t>0$,
\begin{equation}\label{eq:sm-mixed-tail}
 \nu\{(r c_i)^{-1}>t\}
                       \le K_\times t^{-4/3}.
\end{equation}
Consequently, for every Borel $0\le h\le1$ with
$a=\int h\,d\nu$,
\begin{equation}\label{eq:sm-mixed-functional}
 \int\frac{h}{r c_i}\,d\nu
                    \le4K_\times^{3/4}a^{1/4}.
\end{equation}
\end{lemma}
\begin{proof}
Choose any $b>0$.  The part with $r<b$ has mass at most $A_0b^2$.
On the annulus $2^j b\le r<2^{j+1}b$, the condition in
\eqref{eq:sm-mixed-tail} implies
$c_i<(t2^j b)^{-1}$.  Lemma~\ref{lem:sm-intersection}, valid
for every positive value of both arguments, therefore bounds
this annulus by
\[
 A_1(2^{j+1}b)(t2^jb)^{-2}
                     =\frac{2A_1}{t^2b}\,2^{-j}.
\]
The annuli with $j\ge0$ cover all $r\ge b$; most are empty
when $2^jb>\sqrt2$, but keeping them only enlarges the upper sum.
Thus the complete source satisfies
\[
 \nu\{(r c_i)^{-1}>t\}
       \le A_0b^2+\frac{4A_1}{t^2b}.
\]
Take $b=(A_1/A_0)^{1/3}t^{-2/3}$.  Both terms are defined for
all $t>0$, and their sum is
$5A_0^{1/3}A_1^{2/3}t^{-4/3}$.
This choice is a substitution into a finite-count tail bound, not
a collision-dependent residence or a changed physical event.

Layer cake now gives
\[
 \int\frac{h}{r c_i}\,d\nu
 \le\int_0^\infty\min\{a,K_\times t^{-4/3}\}\,dt
                 =4K_\times^{3/4}a^{1/4}.
\]
For $a=0$ use monotone truncation; for $a>0$ split the integral
at $(K_\times/a)^{3/4}$.  The calculation also proves the finite
first moment before it is used anywhere else.
\end{proof}

\begin{theorem}[The full original singular multiplier is uniformly integrable]
\label{thm:sm-full-multiplier}
On the complete regular physical source, with no positive lower
incidence or clearance margin, one has
\begin{equation}\label{eq:sm-full-functional}
 \int h\frac{I_m}{|\nabla_qF|}\,d\nu
       \le J_\times\left(\int h\,d\nu\right)^{1/4}
       \qquad(0\le h\le1).
\end{equation}
In particular, the actual multiplier has the uniform weak bound
\begin{equation}\label{eq:sm-full-tail}
 \nu\{Z_m>t\}\le(J_\times/t)^{4/3}\quad(t>0),
 \qquad \int Z_m\,d\nu\le J_\times.
\end{equation}
For every $1<p<4/3$ it also satisfies
\begin{equation}\label{eq:sm-full-moment}
 \int Z_m^p\,d\nu
                     \le\frac{4}{4-3p}J_\times^p.
\end{equation}
All statements remain valid on the original conditional measure
with the explicitly normalized form
\begin{equation}\label{eq:sm-conditioned-functional}
 \int hZ_m\,d\nu_R^*
       \le J_\times c_*^{-3/4}
                    \left(\int h\,d\nu_R^*\right)^{1/4}
                    \qquad(0\le h\le1).
\end{equation}
Any original event and every original roof test may be retained in
$h$; its values need not factor over time.
\end{theorem}
\begin{proof}
The first inequality of \eqref{eq:sm-intersection} gives
$\nu\{r^{-1}>t\}\le A_0t^{-2}$ for all $t>0$.  Thus
\[
 \int h/r\,d\nu
       \le\int_0^\infty\min\{a,A_0t^{-2}\}\,dt
       =2\sqrt{A_0}\,a^{1/2}
       \le2\sqrt{A_0}\,a^{1/4},
 \qquad a=\int h\,d\nu\le1.
\]
Expand the unchanged $I_m$, use
\eqref{eq:sm-mixed-functional} at every actual interior index,
and sum the deterministic weights
$\sum_{i=1}^{m-1}q_0^{d_i}\le47/3$.
There is no factor $m$, no independent collision law, and no
marked-word pressure estimate in this calculation.  Finally
$|\nabla_qF|=Rr$ and $R\ge R_-$ give exactly
$J_\times$ in \eqref{eq:sm-constants}.  The empty interior sum
for $m=1$ is included.

For $h=\one_{\{Z_m>t\}}$, if its mass $a$ is positive,
$ta\le J_\times a^{1/4}$ proves \eqref{eq:sm-full-tail};
if $a=0$ it is immediate.  For $h=1$ the first-moment estimate
follows directly from \eqref{eq:sm-full-functional}.
Integrating the distribution bound
$\min\{1,(J_\times/t)^{4/3}\}$, split at $t=J_\times$,
proves \eqref{eq:sm-full-moment}.
For the conditional statement apply \eqref{eq:sm-full-functional}
to $\one_{Y_R^*}h$ under $\nu$.  Its integral is
$c_*\int h\,d\nu_R^*$; division of the weighted integral by
$c_*$ leaves exactly $c_*^{-3/4}$.
The kernel proves $\nu\{r=0\}=0$ and invariance proves the
finite-record singular-null assertions.  All operations can first
be performed with bounded multipliers and then passed monotonically
to the actual one; no divergent integral was used to justify them.
\end{proof}

\subsection{The actual insertion on complete critical-origin bands}
The same intersection has a sharper consequence when a physical
critical-origin band is already known to lie in a momentum disk.
It extends the full Borel band estimate of
Theorem~\ref{thm:rc-band} to the actual $I_m$ insertion.

\begin{lemma}[Mixed incidence mass in a momentum disk]
\label{lem:sm-disk}
For every $a>0$ and each $1\le i<m$,
\begin{align}\label{eq:sm-disk}
 \int_{\{r<a\}}c_i^{-1}\,d\nu
             &\le2\sqrt{A_0A_1}\,a^{3/2},\notag\\
 \int_{\{r<a\}}(r c_i)^{-1}\,d\nu
             &\le6\sqrt{A_0A_1}\,a^{1/2}.
\end{align}
Consequently
\begin{equation}\label{eq:sm-disk-Z}
 \int_{\{r<a\}}Z_m\,d\nu
  \le\frac1{R_-}
             \left(2A_0a+94\sqrt{A_0A_1}\,a^{1/2}\right).
\end{equation}
\end{lemma}
\begin{proof}
At every $t>0$, \eqref{eq:sm-intersection} bounds the mass in
the layer-cake integral by
$\min\{A_0a^2,A_1a t^{-2}\}$.
Its integral is $2\sqrt{A_0A_1}\,a^{3/2}$, proving the first
line of \eqref{eq:sm-disk}.  Write its left side as $F_i(a)$.
For $0<r<a$ the identity
$r^{-1}=a^{-1}+\int_r^a s^{-2}\,ds$ and Tonelli give
\[
 \int_{\{r<a\}}\frac1{r c_i}\,d\nu
       =F_i(a)/a+\int_0^a F_i(s)s^{-2}\,ds
       \le6\sqrt{A_0A_1}\,a^{1/2}.
\]
The possible choice of strict or non-strict disk boundaries does
not matter: the joint endpoint kernel gives zero mass to every
circle.  The same formula with $\nu\{r<a\}\le A_0a^2$ gives
$\int_{r<a}r^{-1}\,d\nu\le2A_0a$.
Sum the actual incidence weights, using $6(47/3)=94$, and
divide by $R\ge R_-$ to prove \eqref{eq:sm-disk-Z}.
\end{proof}

\begin{corollary}[Full physical critical-band singular payment]
\label{cor:sm-critical-band}
Let $E_s$ be any measurable one-copy family of actual critical-origin
arm points in Theorem~\ref{thm:rc-band}, with the same $0<\rho<1$
and $A_s=\min\{\rho,\sqrt{2U_m s}\}$.
For every original exact event and Borel $0\le h\le1$,
\begin{equation}\label{eq:sm-critical-band}
 \int_{E_s}\one_\Pi h\frac{I_m}{|\nabla_qF|}\,d\nu_R^*
 \le\frac{2A_0A_s+94\sqrt{A_0A_1}\,A_s^{1/2}}
                 {c_*R_-}.
\end{equation}
In particular the complete actual singular volume is $O(s^{1/4})$
on such bands as $s\downarrow0$, uniformly in all collision counts,
physical radii, actual arms and exact labels.  No tangent-sector
assumption or lower incidence/clearance margin is added.
\end{corollary}
\begin{proof}
The proved physical band geometry gives $E_s\subset\{r\le A_s\}$.
Bound its selector by one and use \eqref{eq:sm-disk-Z}, with the
single original conditioning factor $c_*^{-1}$.  No arm is extended
beyond its actual physical boundary.  As $U_m\le320/9$, the two
terms are respectively $O(s^{1/2})$ and $O(s^{1/4})$.
\end{proof}

\subsection{Full original-roof singular density and deep-source payment}
Use the complete source measures $\mu_{E,H}$ and densities
$f_{E,H}$ from Section~\ref{sec:fr-full-roof}, without any
factor $\one_{\mathcal G_\varepsilon}$.  Define the actual
singular-roof measure
\begin{equation}\label{eq:sm-roof-measure}
 \mu^Z_{E,H}=(L_{m,R})_*
       (Z_m\one_E\one_{\mathcal C_{m,R}}H\,d\nu_R^*).
\end{equation}
The cap and exact-event conventions are unchanged.

\begin{theorem}[Full singular-roof regularity and normalized exceptional volume]
\label{thm:sm-roof}
Whenever $M_E,W_E$ are the full-source bounds of
Lemma~\ref{lem:fr-exhaustion}, the measure
\eqref{eq:sm-roof-measure} is absolutely continuous with density
$f^Z_{E,H}$ and satisfies, for every Borel roof set $S$,
\begin{align}\label{eq:sm-roof}
 \mu^Z_{E,H}(S)
       &\le J_\times c_*^{-3/4}\mu_{E,H}(S)^{1/4},\notag\\
 \|f^Z_{E,H}\|_1
       &\le J_\times c_*^{-3/4}M_E^{1/4},\notag\\
 \|f^Z_{E,H}\|_{12/11,\infty}
       &\le J_\times c_*^{-3/4}(3W_E)^{1/4}.
\end{align}
For $E=M_R$ these are collision-uniform bounds for the entire
actual singular source.  For the complete exceptional union
$E=\mathcal E^{\rm obs}_{D_m,R,m}$, $D_m=\lceil\theta m\rceil$,
$0<\theta<1/2$, they imply uniformly in $R$ and all original
labels and selectors
\begin{align}\label{eq:sm-roof-rates}
 \|f^Z_{D_m,H}\|_1&=O(e^{-\lambda_bm/4}),\notag\\
 \|f^Z_{D_m,H}\|_{12/11,\infty}
                       &=O(e^{-\lambda_bm/12}),\notag\\
 \|f^Z_{D_m,H}\|_p
                       &=O(e^{-\lambda_bm(2/p-7/4)}),
                                 \qquad1<p<12/11.
\end{align}
Every displayed norm has zero limit after multiplication by $m^2$.
In particular the full positive singular volume of this source,
and its volume in every original roof interval, receive an
unconditional normalized payment without a remaining guard class.
\end{theorem}
\begin{proof}
Apply \eqref{eq:sm-conditioned-functional} to
$h=\one_E\one_{\mathcal C_{m,R}}H\one_{\{L_{m,R}\in S\}}$.
This proves the first bound.  Since $\mu_{E,H}$ is finite and
absolutely continuous, the same is true of $\mu^Z_{E,H}$.
In particular a bulk-null critical fiber creates no atom of this
weighted pushforward.  It is a different matter to estimate a
separately defined paired endpoint measure.

For $|S|<\infty$, \eqref{eq:fr-all-borel} gives
\[
 \mu^Z_{E,H}(S)
       \le J_\times c_*^{-3/4}(3W_E)^{1/4}|S|^{1/12}.
\]
Use the superlevel set of the nonnegative density to obtain the
weak $L^{12/11}$ bound, exactly as in the proof of
\eqref{eq:fr-weighted-full}.  The total mass follows by setting
$S=\mathbb R$ in the first inequality.
For the exceptional source, Theorem~\ref{thm:fr-deep-source}
gives $M_{D_m}=O(e^{-\lambda_bm})$ and
$W_{D_m}=O(e^{-\lambda_bm/3})$.  Apply the indicated fourth
powers and then Lemma~\ref{lem:fr-strong} with $q=12/11$.
The exponent in the strong $p$th power is
$(12-11p)/4+(p-1)=2-7p/4$, as stated.
The exponent $2/p-7/4$ is positive throughout the indicated
range.  The total-volume bound controls every roof interval
without changing its width or location.
\end{proof}

\begin{corollary}[The complete regular-interior variation budget]
\label{cor:sm-variation}
Use exactly the retained countable one-copy physical sheet space of
Proposition~\ref{prop:el-positive}, with all original exact-event
and section cuts.  For an original guard parameter $\varepsilon>0$,
a Borel physical source $E$ and a physical selector $0\le H\le1$,
define the nonnegative roof measure
\[
 \mathcal V_{\varepsilon,E,H}(S)
  =\int\one_S(u)\one_E(\Psi(\theta,u))H(\Psi(\theta,u))
          |\partial_u f^\varepsilon(u,\theta)|\,du\,d\theta.
\]
The derivative is taken only on the regular interior of each actual
sheet, before multiplying by any indicator.  Then
\begin{equation}\label{eq:sm-variation}
 \mathcal V_{\varepsilon,E,H}(S)
 \le C_{\rm reg}(1+\varepsilon^{-1})
                       \mu^Z_{E,\one_\Pi H}(S)
 \qquad\hbox{for every Borel roof set }S.
\end{equation}
In particular the complete regular-interior variation is finite,
uniformly in collision count and radius, with total bound
\[
 C_{\rm reg}(1+\varepsilon^{-1})J_\times c_*^{-3/4}.
\]
Its density has the corresponding full weak $L^{12/11}$ and
strong $L^p$, $1<p<12/11$, bounds.  On
$E=\mathcal E^{\rm obs}_{D_m,R,m}$, all the rates in
\eqref{eq:sm-roof-rates} hold for this variation density, multiplied
by the fixed factor $C_{\rm reg}(1+\varepsilon^{-1})$.
Consequently its $m^2$-normalized total variation and all those
roof norms tend to zero.  In the original application
$\varepsilon=\varepsilon(B)$ is fixed before the collision limit.
\end{corollary}
\begin{proof}
Apply Proposition~\ref{prop:el-positive} to the nonnegative
pullback selector
$\one_S(u)\one_E(\Psi(\theta,u))H(\Psi(\theta,u))$.
The retained physical one-copy coarea identifies its upper side
with precisely the right side of \eqref{eq:sm-variation}.
The exact event is unchanged, and its endpoint tests put it inside
$\mathcal C_{m,R}$ as required in \eqref{eq:sm-roof-measure}.
Theorem~\ref{thm:sm-roof} proves finiteness and the asserted
norms.  Positivity and Radon--Nikodym differentiation of the
measure inequality transfer its density bounds to the variation
measure.  For the deep source, apply \eqref{eq:sm-roof-rates}
after the complete physical sum.  Neither $H$ nor an exact
section or roof indicator is differentiated.

This argument uses exactly the inherited one-copy coarea and the
already native regular-variation inequality; it does not reconstruct
a new pairing.  On a whole sheet the existing absolute-continuity
and trace result identifies the interior derivative integral with
the ordinary variation of $f^\varepsilon$ on its closed physical
segment.  The present localized statement is the restriction of
that derivative measure, not the derivative of an indicator-cut
function.  It therefore does not count jumps between different
segments, manufacture a singular trace charge, or bound a positive
endpoint minimum by an unrelated small bulk set.
\end{proof}

\begin{corollary}[The actual singular insertion in the full marked comparison]
\label{cor:sm-marked}
Use the exact full-source marked comparison of
Corollary~\ref{cor:fr-full-comparison}.  With its actual additional
multiplier $Z_m=I_m/|\nabla_qF|$, one has
\begin{equation}\label{eq:sm-marked}
 \left\|d\mathfrak D_H^Z/dt\right\|_p
   \le\|f^Z_{d,1}\|_p+
     2R_+|s|q_{\rm a}^{\lfloor d/N\rfloor}\|f^Z_1\|_p,
                  \qquad1\le p<12/11.
\end{equation}
Its $m^2$-normalized value tends to zero uniformly in all lattice
and return twists and $|s|\le e^{\lambda m}$, for every fixed
$0<\lambda<\lambda_c$.  Every collision and the signed endpoint
clock work are retained, with the same negative compensating phase.
\end{corollary}
\begin{proof}
Multiply the actual pointwise phase domination in the proof of
Corollary~\ref{cor:fr-full-comparison} by the nonnegative $Z_m$.
Theorem~\ref{thm:sm-roof} and Lemma~\ref{lem:fr-strong} give
the required full-source strong bounds and exceptional-source
decay.  The remaining clock factor decays at rate
$\lambda_c-\lambda>0$.  No Fourier integration over an
unbounded range is asserted.
\end{proof}

The physical significance of \eqref{eq:sm-full-functional} is that
the two possible singular mechanisms, an interior grazing contact
and a vanishing full-roof gradient, have a quantified joint bound
on the real orbit.  Separate marginal estimates would give only
a borderline product and would not prove this first moment.
The intersection of two already valid physical block kernels is
the additional mechanism.  It completes the full bulk singular
integrability and the complete deep-source singular-volume payment.
It does not bound a positive short-segment endpoint trace by that
volume without its own proved charge inequality.  Nor does weak
$L^{12/11}$ roof control establish a pointwise roof supremum.
The original shallow positive source, full comparison capacity,
singular pairing interfaces and the final spectral, arithmetic and
conditional pointwise inversion requirements retain their stated
scope and the original order of limits.
